\section{Assumptions}
\label{sec:assumptions}
In this work, we categorize the assumptions into two groups: those necessary for 
the identification of causal effects and those required for estimation. 
As estimation assumptions are method-specific and highly technical, 
we refer the reader to the original papers for further detail. 
In this section, we outline the common \emph{identification} assumptions 
behind both the FE--MSM and SeqGPLVM approaches. Both methods described in this section 
rely on a set of assumptions to ensure the identifiability
of the causal effects in the presence of unmeasured time-invariant confounding. 
We summarize these assumptions below:
% --- Common assumptions (C1, C2, ...) ---
% Assumes you already defined the assumpC environment (C-prefix) in the preamble.

\begin{assumpC}[Consistency]\label{ass:C-consistency}
The observed outcome for each individual under their observed treatment history equals
the potential outcome under that treatment history. Formally, for each individual $i$
and time $t$, if $\bar{A}_{i,t}=\bar{a}_{i,t}$, then $Y_{i,t+1}=Y_{i,t+1}(\bar{a}_{i,t})$. 
Moreover, $Y_{i,T+1} = Y_{i,T+1}(\underline{a}_{i,T-k})$ if 
$\bar{A}_{i,T}=\left(\bar{A}_{i,T-k-1},\underline{a}_{i,T-k}\right)$.
\end{assumpC}

\begin{assumpC}[Positivity]\label{ass:C-positivity}
\[
\Pr\!\left[ A_t = a_t \,\middle|\, 
\bar{A}_{t-1} = \bar{a}_{t-1},\,
\bar{X}_t = \bar{x}_t \right] > 0
\]

for all values $(\bar{a}_{t-1}, \bar{x}_t)$ with

\[
\Pr\!\left[ \bar{A}_{t-1} = \bar{a}_{t-1},\,
\bar{X}_t = \bar{x}_t \right] \neq 0
\]

in the population of interest, and for all $a_t \in \mathcal{A}_t$.
\end{assumpC}

\begin{assumpC}[No interference]\label{ass:C-no-interference}
The treatment assignment for one individual does not affect the potential outcomes of
other individuals.
\end{assumpC}

\begin{assumpC}[Time-invariance of unmeasured confounding]\label{ass:C-timeinvariant-uc}
The unmeasured confounding is assumed to be
time-invariant. That is, $U_{i,t} = U_i$ for $t \in \{1,\dots,T\}$. We can then write 
conditional ignorability as 
\[
    \{ Y_{i,t+1}(\bar{a}_{t}),\, \bar X_{i,t+1}(\bar a_t) \}
\;\perp\!\!\!\perp\;
A_{it}
\;\big|\;
\bar X_{it},\,
\bar A_{i,t-1} = \bar a_{t-1},\,
U_i. \]
\end{assumpC}

